%  04_methods.tex
\section{Methods}
\label{sec:methods}

We compare two families of face recognition models distinguished by their training
objective: \emph{pair-based} models that learn directly from same/different pairs, and
\emph{embedding-based} models that learn a compact hyperspherical representation via
closed-set classification with an angular margin.  Both families share a common inference
interface: given two images, cosine similarity of L2-normalised embeddings is computed and
compared to a learned threshold.

\subsection{Pair-based: Siamese Convolutional Network}

The Siamese network~\cite{bromley1994signature, koch2015siamese} consists of two weight-tied
embedding towers and a pairwise decision head.

\paragraph{Embedding tower.}
The tower processes $105\!\times\!105$ RGB images through four convolutional stages with
``valid'' (zero-padding-free) convolutions:
\begin{enumerate}[leftmargin=1.5em, itemsep=1pt, parsep=0pt]
  \item Conv $64 @ 10{\times}10$, ReLU $\rightarrow$ MaxPool $2{\times}2$ (stride 2):
        $105\!\to\!96\!\to\!48$
  \item Conv $128 @ 7{\times}7$, ReLU $\rightarrow$ MaxPool $2{\times}2$:
        $48\!\to\!42\!\to\!21$
  \item Conv $128 @ 4{\times}4$, ReLU $\rightarrow$ MaxPool $2{\times}2$:
        $21\!\to\!18\!\to\!9$
  \item Conv $256 @ 4{\times}4$, ReLU: $9\!\to\!6$
  \item Flatten $(256 \cdot 6 \cdot 6 = 9\,216)$ $\rightarrow$ Linear$(9\,216, 4\,096)$ $\rightarrow$ Sigmoid
\end{enumerate}
Producing a 4\,096-dimensional embedding with values in $[0, 1]$.

\paragraph{Pair head.}
The absolute L1 distance between the two embeddings, $|\mathbf{e}_a - \mathbf{e}_b|$, is
passed through a final linear layer (4\,096~$\to$~1) followed by a sigmoid, producing a
similarity score $s \in [0, 1]$.  L2 regularisation with weight $10^{-4}$ is applied via
the optimiser, replicating the Keras \texttt{l2(1e-4)} behaviour.

\paragraph{Loss.}
Siamese networks are trained with the Focal
loss~\cite{lin2017focal}: $\mathrm{FL}(p_t) = -\alpha_t (1-p_t)^\gamma \log p_t$,
with $\alpha = 0.75$, $\gamma = 2.0$.  Focal loss down-weights easy negatives that
dominate pairwise training sets, improving calibration of the decision boundary.

\paragraph{Inference.}
At inference time only the embedding tower is used; the trained pair-head is discarded.
Verification scores are computed as the cosine similarity of L2-normalised embeddings,
providing a consistent metric space across all three model families.

Figure~\ref{fig:architectures} (left) illustrates the Siamese architecture.

\subsection{Embedding-based: ArcFace and AdaFace}

Both embedding models share a common backbone and projection block, differing only in their
classification head.  The architecture maps $112\!\times\!112$ RGB images to a
512-dimensional hyperspherical embedding.

\paragraph{Backbone.}
A ResNet-50~\cite{he2016resnet} initialised with ImageNet1K-V2 weights serves as the
feature extractor.  The original fully connected classification head is replaced by
\texttt{nn.Identity()}, yielding a 2\,048-dimensional feature vector.

\paragraph{Projection block.}
Following standard practice~\cite{deng2019arcface}, the 2\,048-d backbone output is
compressed by:
$\mathrm{BN}(2048) \!\to\! \mathrm{Dropout}(0.0) \!\to\! \mathrm{Linear}(2048, 512,
\text{bias=False}) \!\to\! \mathrm{BN}(512)$.
This produces a 512-d embedding $\mathbf{z} \in \mathbb{R}^{512}$ (unnormalised).

\paragraph{ArcFace head~\cite{deng2019arcface}.}
Let $\hat{\mathbf{z}} = \mathbf{z}/\|\mathbf{z}\|$ be the L2-normalised embedding and
$\hat{\mathbf{w}}_k$ the L2-normalised class weight for identity $k$.  Define the angle
$\theta_{y}$ between the embedding and the ground-truth class weight.  The ArcFace logit is:
\begin{equation}
  \ell_y = s \cdot \cos(\theta_y + m), \qquad \ell_{k \ne y} = s \cdot \cos \theta_k
  \label{eq:arcface}
\end{equation}
where $m = 0.2$ (margin, radians) and $s = 30$ (scale) in the tuned configuration used in
our experiments; logits are passed through cross-entropy.

\paragraph{AdaFace head~\cite{kim2022adaface}.}
AdaFace augments the ArcFace formulation with an image-quality proxy derived from
$\|\mathbf{z}\|$, the feature norm of the unnormalised embedding.
A per-batch running mean $\mu_b$ and standard deviation $\sigma_b$ of feature norms are
maintained via exponential moving average (EMA momentum $\alpha = 0.01$).
The margin is then adapted per sample:
\begin{align}
  g_{\text{margin}} &= \operatorname{clip}\!\left(\frac{\|\mathbf{z}\| - \mu_b}
                                                {\sigma_b + \varepsilon} \cdot h,\,
                                                -1, 1\right), \\
  \ell_y &= s \cdot \bigl[\cos(\theta_y - m \cdot g_{\text{margin}}) - m \cdot
                          g_{\text{margin}}\bigr]
  \label{eq:adaface}
\end{align}
with $h = 0.333$, $m = 0.2$, $s = 30$.  Images with below-average feature norms (low
quality) receive a reduced margin, lessening their influence on the decision boundary.

\paragraph{Inference.}
At test time both ArcFace and AdaFace return the L2-normalised embedding
$\hat{\mathbf{z}} = \mathbf{z}/\|\mathbf{z}\|$ without the classification head.
Verification uses cosine similarity as in the Siamese case.

Figure~\ref{fig:architectures} summarises both training topologies.

\begin{figure}[t]
\centering
% Two sub-figures side by side, each with its own tikzpicture
\begin{subfigure}[t]{0.46\linewidth}
\centering
\begin{tikzpicture}[node distance=0.42cm]
  % --- Siamese pair topology ---
  \node[block, fill=colSiamese!10, draw=colSiamese!50,
        minimum width=1.45cm, font=\small] (imgA) {Image $A$};
  \node[block, fill=colSiamese!10, draw=colSiamese!50,
        minimum width=1.45cm, font=\small,
        right=1.0cm of imgA] (imgB) {Image $B$};

  \node[block, fill=colSiamese!20, draw=colSiamese!60,
        minimum width=1.45cm, minimum height=1.1cm,
        align=center, font=\scriptsize,
        below=of imgA] (towerA) {Conv\\Tower};
  \node[block, fill=colSiamese!20, draw=colSiamese!60,
        minimum width=1.45cm, minimum height=1.1cm,
        align=center, font=\scriptsize,
        below=of imgB] (towerB) {Conv\\Tower};

  % tie arrow between towers
  \draw[<->, thick, dashed, colSiamese!60]
    ([yshift=2pt]towerA.north east) -- ([yshift=2pt]towerB.north west)
    node[midway, above, font=\tiny, colSiamese]{shared};

  \node[block, fill=colSiamese!15, draw=colSiamese!50,
        minimum width=1.45cm, font=\scriptsize,
        below=of towerA] (embA) {$\mathbf{e}_A$};
  \node[block, fill=colSiamese!15, draw=colSiamese!50,
        minimum width=1.45cm, font=\scriptsize,
        below=of towerB] (embB) {$\mathbf{e}_B$};

  % L1 node: centred horizontally between the two embedding nodes
  \coordinate (midEmb) at ($(embA.south)!0.5!(embB.south)$);
  \node[block, draw=gray!50, font=\scriptsize,
        minimum width=3.2cm, below=0.7cm of midEmb] (l1)
    {$|\mathbf{e}_A - \mathbf{e}_B|$};
  \node[block, draw=gray!50, font=\scriptsize,
        minimum width=3.2cm, below=0.4cm of l1] (head)
    {Linear$(4096,1)$ + $\sigma$};
  \node[block, fill=colArcFace!10, draw=colArcFace!40, font=\scriptsize,
        minimum width=3.2cm, below=0.4cm of head] (loss)
    {Focal Loss};

  \draw[arrow] (imgA) -- (towerA);
  \draw[arrow] (imgB) -- (towerB);
  \draw[arrow] (towerA) -- (embA);
  \draw[arrow] (towerB) -- (embB);
  \draw[arrow] (embA.south) -- ++(0,-0.35) -| (l1.north west);
  \draw[arrow] (embB.south) -- ++(0,-0.35) -| (l1.north east);
  \draw[arrow] (l1)   -- (head);
  \draw[arrow] (head) -- (loss);
\end{tikzpicture}
\caption{Pair family (Siamese)}
\end{subfigure}
\hfill
\begin{subfigure}[t]{0.46\linewidth}
\centering
\begin{tikzpicture}[node distance=0.42cm]
  % --- Embedding topology ---
  \node[block, fill=colArcFace!10, draw=colArcFace!50,
        minimum width=3.2cm, font=\small] (img)
    {Image $\mathbf{x}$};
  \node[block, fill=colArcFace!15, draw=colArcFace!50,
        minimum width=3.2cm, minimum height=0.75cm,
        font=\scriptsize, below=of img] (bb)
    {ResNet-50 backbone};
  \node[block, draw=gray!50, minimum width=3.2cm, font=\scriptsize,
        below=of bb] (proj)
    {BN $\!\to\!$ Drop $\!\to\!$ FC(512) $\!\to\!$ BN};
  \node[block, fill=colArcFace!15, draw=colArcFace!50,
        minimum width=3.2cm, font=\scriptsize,
        below=of proj] (emb)
    {$\mathbf{z}\in\mathbb{R}^{512}$};
  \node[block, fill=colArcFace!20, draw=colArcFace!60,
        minimum width=3.2cm, minimum height=0.9cm,
        align=center, font=\scriptsize, below=of emb] (headE)
    {Angular Margin Head\\(ArcFace or AdaFace)};
  \node[block, fill=colArcFace!10, draw=colArcFace!40,
        minimum width=3.2cm, font=\scriptsize,
        below=of headE] (lossE)
    {Cross-entropy Loss};

  \draw[arrow] (img)   -- (bb);
  \draw[arrow] (bb)    -- (proj);
  \draw[arrow] (proj)  -- (emb);
  \draw[arrow] (emb)   -- (headE);
  \draw[arrow] (headE) -- (lossE);

  % inference side note
  \draw[dashed arrow] (emb.east) -- ++(0.9,0)
    node[right, font=\tiny, align=left]{infer:\\$\hat{\mathbf{z}}$};
  \node[font=\tiny, gray, below=0.1cm of lossE]
    {(head removed at inference)};
\end{tikzpicture}
\caption{Embedding family (ArcFace / AdaFace)}
\end{subfigure}
\caption{Training topologies for the two model families.
\emph{Left}: the Siamese pair-based network uses weight-tied towers and an L1+sigmoid
decision head trained with Focal loss.
\emph{Right}: the embedding-based network uses a ResNet-50 backbone with a projection
block and an angular-margin head.  At inference the head is discarded and cosine
similarity on L2-normalised embeddings drives verification.}
\label{fig:architectures}
\end{figure}

\subsection{Hyperparameter tuning}

Table~\ref{tab:hparams} summarises the final training configurations.
The embedding-model default configuration (SGD, $\eta=0.1$, $m=0.5$, $s=64$) collapsed
completely on our small datasets (ROC-AUC $= 0.50$, Section~\ref{sec:experiments}), as the
paper-scale learning rate destroyed the ImageNet pretrained backbone before meaningful
features were learned.  The tuned configuration (Adam, $\eta=10^{-3}$, $m=0.2$, $s=30$)
was selected by Optuna~\cite{akiba2019optuna} hyperparameter search on the validation split.

\begin{table}[h]
\centering
\caption{Training hyperparameters.}
\label{tab:hparams}
\setlength{\tabcolsep}{5pt}
\begin{tabular}{lll}
\toprule
Hyperparameter & Siamese & ArcFace / AdaFace \\
\midrule
Optimiser        & Adam           & Adam \\
Learning rate    & $5\times10^{-5}$ & $1\times10^{-3}$ \\
Weight decay     & $0.0$          & $5\times10^{-4}$ \\
LR schedule      & constant       & cosine + 1 warmup epoch \\
Max epochs       & 80             & 40 \\
Early-stop monitor & \recall      & val ROC-AUC \\
Early-stop patience & 5           & 7 \\
Mixed precision  & bf16           & bf16 \\
EMA decay        & 0.0            & 0.999 \\
Loss margin ($m$)& ---            & 0.2 \\
Loss scale ($s$) & ---            & 30 \\
Focal $\alpha / \gamma$ & 0.75 / 2.0 & --- \\
\bottomrule
\end{tabular}
\end{table}
